%% Para compilar o programa, use pdflatex
%% Opcoes:
%%
%% Cursos:
%% cic: Bacharelado em Ciência da Computação
%% tsi: Técnologo em Sistemas para Internet
\documentclass[tsi]{ifbclass/ifbclass}

\address{BRASÍLIA}

\title{Trabalho de Conclusão de Curso}

\date{2019}

\author{Marcos Vasconcellos de Andrade}
\adviser{Alisson Silva}

\membroum[a]{Dr. Daniel Lima}
\membrodois{Terceiro Membro de Banca}
%\membrotres[a]{Dr.ª Terceira Membro da Banca}
%\membroquatro[a]{Dr.ª Quarta Membro da Banca}

%\coadviser{Nome completo do co-orientador (se existir)} %TODO: nao suportado

% Macros (Se necessario, defina suas proprias aqui)
\def\x{\checkmark}
%\let\lstlistoflistings\origlstoflistings
  \begin{document}

\frontmatter

\frontpage

\presentationpage

% Quando a biblioteca preparar a ficha catalografica, insira-a aqui
% Antes da defesa do trabalho, crie uma ficha falsa
\begin{fichacatalografica}
  \FakeFichaCatalografica
%     \includepdf{fig_ficha_catalografica.pdf} % crie o arquivo come esse nome e descomente essa linha
\end{fichacatalografica}

\banca

%Comente/Remova todo o conteudo entre \begin{dedicatory} e \end{dedicatory} para remover.
\begin{dedicatory} %OPCIONAL
Dedico este trabalho à minha esposa Rosimery que me apoiou e compreendeu as
várias horas empenhadas para a conclusão deste curso, inclusive aos finais
de semana e feriados.
\end{dedicatory}
  
\acknowledgements %OPCIONAL Comente a linha abaixo para remover.
Agradeço ao meu orientador Prof. Dr. Alisson Silva, que aceitou em ser meu
orientador logo no primeiro contato...

Aos meus colegas de sala, pelo companherismo e apoio, 

A Secretaria do Curso, pela cooperação.

Gostaria de deixar registrado também, o meu reconhecimento à minha família, pois acredito que sem o apoio deles seria muito difícil vencer esse desafio. 

Enfim, a todos os que por algum motivo contribuíram para a realização desta pesquisa.


%Novamente, comente todo conteudo entre {epigraph} para remover
\begin{epigraph}[]{Nome do autor} %OPCIONAL
Elemento opcional.

Espaço destinado à epígrafe (elemento opcional). Nesta folha, o autor usa uma
citação, seguida de indicação de autoria e ano, relacionada com a matéria
tratada no corpo do trabalho.
\end{epigraph}

\resumo
% Escreva seu resumo no arquivo resumo.tex
{\parindent0pt
  VASCONCELLOS DE ANDRADE, Marcos. Sistema de Administração de Condomínio Personalizado.  2019. 65 f. 
Trabalho de Conclusão de Curso (Graduação) – Tecnólogo em Sistemas para Internet. 
Instituto Federal de Brasília – Campus Brasília. Brasília/DF, 2019.
\vspace{1cm}

Elemento obrigatório, constituído de uma sequência de frases concisas e objetivas,
fornecendo uma visão rápida e clara do conteúdo do estudo. O texto deverá conter no
máximo 500 palavras e ser antecedido pela referência do estudo, com exceção do resumo
inserido no próprio documento. Também, não deve conter citações. O resumo deve ser redigido
em parágrafo único, espaçamento simples e seguido das palavras representativas do conteúdo
do estudo, isto é, palavras-chave, em número de três a cinco, separadas entre si por ponto e
finalizadas também por ponto. Usar o verbo na terceira pessoa do singular, com linguagem
impessoal (pronome SE), bem como fazer uso, preferencialmente, da voz ativa.


\begin{keywords}
%Primeira palavra. Segunda palavra. Terceira palavra. Quarta palavra. Quinta-palavra.
condomínio. síndico. condômino. 
\end{keywords}

}
  
\abstract
% Escreva seu abstract em ingles no arquivo abstract.tex
{\parindent0pt
  SOBRENOME, Prenome do Autor do Trabalho. Título do trabalho: subtítulo (se houver).  2018. 65 f. 
Trabalho de Conclusão de Curso (Graduação) – Tecnólogo em Sistemas para Internet. 
Instituto Federal de Brasília – Campus Brasília. Brasília/DF, 2018.
\vspace{1cm}

Elemento obrigatório, constituído de uma sequência de frases concisas e objetivas,
fornecendo uma visão rápida e clara do conteúdo do estudo. O texto deverá conter no
máximo 500 palavras e ser antecedido pela referência do estudo, com exceção do resumo
inserido no próprio documento. Também, não deve conter citações. O resumo deve ser redigido
em parágrafo único, espaçamento simples e seguido das palavras representativas do conteúdo
do estudo, isto é, palavras-chave, em número de três a cinco, separadas entre si por ponto e
finalizadas também por ponto. Usar o verbo na terceira pessoa do singular, com linguagem
impessoal (pronome SE), bem como fazer uso, preferencialmente, da voz ativa.

\begin{keywords}
Keyword. Second keyword. Third keyword. Keyword.
\end{keywords}

}

%ATENCAO.
%Se alguma dessas listas estiver vazia no seu trabalho, comente a linha com um '%'
% Lista de figuras
\listoffigures

% Lista de algortimos
\lstlistoflistings

% Lista de tabelas
\listoftables

% Lista de acronimos
% Acronyms manual: http://linorg.usp.br/CTAN/macros/latex/contrib/acronym/acronym.pdf
\listofacronyms
\begin{acronym}[ACRONYM] 
% Change the word ACRONYM above to change the acronym column width.
% The column width is equals to the width of the word that you put.
% Read the manual about acronym package for more examples:
%   http://linorg.usp.br/CTAN/macros/latex/contrib/acronym/acronym.pdf
\acro{afm}[AFM]{Alphabet Frequency Matrix}
\acro{api}[API]{Application Programming Interface}
\acro{arima}[ARIMA]{Auto-Regressive Integrated Moving Average}
\acro{brn}[BRN]{Bug Report Network}
\acro{bts}[BTS]{Bug Triage System}
\acro{cas}[CAS]{Context-Aware Systems}
\acro{ccb}[CCB]{Change Control Board}
\acro{cr}[CR]{Change Request}
\acro{cvs}[CVS]{Concurrent Version System}
\acro{es}[ES]{Expert System}
\acro{floss}[FLOSS]{Free/Libre Open Source Software}
\acro{glr}[GLR]{Generalized Linear Regression}
\acro{gqm}[GQM]{Goal Question Metric}
\acro{html}[HTML]{HyperText Markup Language}
\acro{ir}[IR]{Information Retrieval}
\acro{irt}[IRT]{Recôncavo Institute of Technology}
\acro{jdt}[JDT]{Jazz Duplicate Finder}
\acro{lda}[LDA]{Latent Dirichlet Allocation}
\acro{loc}[LOC]{Lines of Code}
\acro{lsi}[LSI]{Latent Semantic Indexing}
\acro{ms}[MS]{Mapping Study}
\acro{msr}[MSR]{Mining Software Repositories}
\acro{nlp}[NLP]{Natural Language Processing}
\acro{promise}[PROMISE]{Predictive Models in Software Engineering}
\acro{rbes}[RBES]{Rule-Based Expert System}
\acro{rhel}[RHEL]{RedHat Enterprise Linux}
\acro{saas}[SaaS]{Software as a Service}
\acro{scm}[SCM]{Software Configuration Management}
\acro{serpro}[SERPRO]{Brazilian Federal Organization for Data Processing}
\acro{slr}[SLR]{Stepwise Linear Regression}
\acro{slr}[SLR]{Systematic Literature Review}
\acro{svd}[SVD]{Singular Value Decomposition}
\acro{svm}[SVM]{Support Vector Machine}
\acro{svn}[SVN]{Subversion}
\acro{tfidf}[TF-IDF]{Term Frequency-Inverse Document Frequency}
\acro{vsm}[VSM]{Vector Space Model}
\acro{xp}[XP]{Extreming Programming}
\end{acronym}

% Sumario
\tableofcontents

\mainmatter

%Para cada seção do seu trabalho, edite o arquivo .tex na pasta text/
\chapter{Introdução}
\label{chp:introduction}
%Faça aqui, uma introdução geral da área do conhecimento à qual o tema escolhido
%está ligado. 
O tema está diretamente ligado a área de desenvolvimento de software e terá como
objetivo o desenvolvimento de uma Aplicação Web e Mobile utilizando a linguagem
de programação JavaScript. No frontend será utilizado o React, no backend será
utilizado o Nodejs e o banco de dados será o PostGresql. O React Native
possibilitará a portabilidade do sistema para mobile.
\section{Tema}
% A melhor forma de determinar o tema abordado é através de hipóteses. A hipótese
% consiste em uma afirmativa que você considera verdadeira e que vai provar ou
% buscar provar ao longo de seu trabalho. Outra forma é delimitando o problema em
% forma de uma pergunta de partida. Apresente uma visão geral do assunto que será
% abordado no trabalho.
O tema escolhido será desenvolver um Sistema de Administração de Condomínios
personalizado. Atualmente, existem sistemas com soluções personalizadas e acessíveis
para administração de condomínios?
\section{Problema}
% Dedique este tópico a esclarecer o que o pretende de fato com o seu esforço de
% pesquisa. Problema é a questão a ser respondida pelo trabalho, que motivou a sua
% realização. É uma questão que já tomou se formou em sua mente, derivada de
% teorias da área pesquisada e de sua observação sobre um fenômeno.  Normalmente
% se utilizam os subitens abaixo como meios de se determinar claramente os
% objetivos, o que também colabora para a delimitação do escopo do trabalho. Está
% estreitamente ligado ao objetivo geral, que, normalmente, consiste em encontrar
% a resposta para o problema de pesquisa.  O que você viu que é um problema que
% precisa de solução? É viável? Você consegue fazer? O problema é sempre uma
% dificuldade, uma lacuna.
Por morar em um condomínio, foi observado que o síndico do bloco não possui
nenhum sistema para administração de condomínio e que mensalmente são encadernados
mais de 100(cem) páginas dentre comprovantes, prestações de conta, notas fiscais,
controle de empregados, etc.

Ao ser realizada uma pesquisa nos condomínios adjacentes foi percebido a mesma
problematica e que para resolver essas dores poderia surgir uma excelente
oportunidade de negócio.

Ao ser realizado a entrevista inicial com o síndico do bloco, obteve-se o
aceite de imediato e o este condomínio será o piloto no desenvolvimento
do TCC, que em troca receberá o sistema com custo zero.

\subsection{Objetivo geral}
% É a resposta ao problema especificado acima, ou seja, aquilo que se pretende
% fazer e que, depois de atingido, estará concluído o trabalho.. Alguns verbos
% utilizados para determinar o objetivo geral: contribuir / facilitar / subsidiar
% / propor / clarear / permitir / agregar / compreender.
Desenvolver uma aplicação Web para administração de condomínios e que
terá como objetivo geral a criação de um portal, onde serão oferecidos
os serviços e que encontra-se hospedado em \url{http://cliquesindico.com.br}.
Este portal será um "cartão de visitas" para os futuros contratos de
desenvolvimento de sistemas personalizados para condomínios.

A idéia principal é modularizar a aplicação através de subdomínios, onde cada 
bloco ou condomínio terá um subdomínio único e com um sistema personalizado
de acordo com a necessidade do condomínio.

Para o condomínio que será piloto neste projeto, foi criado o subdomínio
\url{http://sqn306b.cliquesindico.com.br}, onde na página inicial terá uma página
de login, onde inicialmente será definidos o perfil de administrador e de usuário
comum. Após se logar na aplicação haverá um redirecionamento para a página
principal do condomínio, onde encontrará um mural com avisos gerais do condomínio.

Serão enviados e-mail aos condominos com avisos, notificações, cobranças, etc

Será criado uma página para prestação de contas mensal do condomínio.

Após a conclusão do sistema, o síndico e a diretoria do condomínio terão total
autonomia para inserir novos avisos, atualizar a prestação de contas e imprimir
relatórios.

\subsection{Objetivos específicos}
% Os objetivos específicos detalham os objetivos gerais através de etapas ou fases
% de pesquisa. Devem ser utilizados verbos no infinitivo, assinalando as ações
% propostas para alcançar o objetivo geral. Os verbos utilizados aqui são os de
% ação, que serão utilizados na metodologia.
\begin{itemize}
    \item Realizado a entevista inicial para determinação da dores e levantamento
    inicial de requisitos.
    
    \item Criar o cronograma do Projeto.
    
    \item Criação de um portal onde serão inseridas as informações gerais sobre
negócio, que terá como diferencial a personalização e o preço competitivo.
    
    \item Cada condomínio terá um subdomínio personalizado e idependente do
restante do sistema.

    \item Será criado um módulo para login no sistema do condomínio em questão
que terá a opção de cadastro, caso o usuário ainda não esteja cadastrado,
e o acesso será realizado através e-mail e senha. O sistema notificará o síndico
do condomínio correspondente, a fim de que valide o cadastro ou revogue o acesso
caso o usuário não tenha relação com o condomínio.

    \item Será recolhido junto a administração do condomínio a documentação 
inicial para que o sistema começe a ser desenvolvido, tais como, e-mail
dos moradores, avisos iniciais para inserção na página principal e modelo de
ajuste de contas.

    \item Desenvolvimento do sistema.
    
    Desenvolver o frontend, backend, api e banco de dados. Para isso será 
utilizado a linguagem JavaScript, React no frontend, NodeJS para APIs e para
o backend e Postgreesql para o banco de dados.

    Para modelagem do banco será utilizado o BrModelo e para criação do banco o PGAdmin.

    \item Realização dos testes para validação inicial do sistema, junto a administração
do condomínio e levantamento de problemas e melhoramentos.

    \item Documentar o sistema
    
\end{itemize}





Implementar os diferentes perfis de acesso, de acordo com a real necessidade
do condomínio

\section{Estrutura do TCC}
Neste item você vai descrever como está constituída a monografia, indicando o
que será encontrado em cada uma das sessões seguintes.

\subsection{Classificação da Pesquisa}
Neste item será apresentada a classificação da pesquisa quanto aos objetivos
(exploratória, descritiva ou explicativa); aos procedimentos (Pesquisa
bibliográfica, Pesquisa documental, Pesquisa experimental, Estudo de caso
controle, Levantamento, Estudo de caso ou Estudo de campo) e ao método de
investigação científica (qualitativa ou quantitativa).

\chapter{Conceitos gerais e revisão da literatura}
Neste capítulo deve ser proporcionado o estado da arte / referencial teórico
sobre o tema a que se refere o estudo. Um bom pesquisador não deve repetir
trabalhos já concluídos ou que já estão em andamento. Por isso esta sessão é
onde o autor demonstra até onde vai a pesquisa atual no campo de estudos em
questão e estabelece as bases sobre as quais desenvolverá o estudo proposto. A
seguir são mostrados alguns exemplos de como deve-se inserir as figuras e
tabelas. A Figura \ref{fig:exemplo} mostra um exemplo de como inserir uma
figura no texto. A Tabela \ref{tb:exemplo} mostra o exemplo de como uma tabela
deve ser inserida.  Voce pode referenciar capítulos e seções adicionando labels
à elas. Por exemplo, descrevemos a introdução no Capítulo
\ref{chp:introduction}.

\begin{figure}[htb]
    \centering
    \includegraphics[width=0.45\textwidth]{images/figura.png}
    \caption{Exemplo de como inserir Figura}
    \label{fig:exemplo}
\end{figure}
%Altere o numero junto ao 'width' para acertar o tamanho da figura. Coloque
%numeros maior que 0 e até 1

\begin{table}[htb]
\caption{Modelo de como as tabelas devem ser inseridas no texto}
\label{tb:exemplo}
\centering
\begin{tabular}{|l|c|r|r|} %left, center, right. Você pode mudar isso
\hline
Índice  & Coluna 1 & Coluna 2 & Coluna 3 \\
\hline
Linha 1 &          &          &          \\
Linha 2 &          &          &          \\
Linha 3 &          &          &          \\
\hline
\end{tabular}
\end{table}
%Para manter a sanidade, recomenda-se deixar sempre os '&' alinhados em todas
%as linhas e colunas

\chapter{Metodologia}

Aqui conterão os métodos e procedimentos adotados no desenvolvimento do
trabalho. Esta é uma das sessões mais importantes pois demonstra o poder
científico que foi utilizado para a pesquisa. Sem uma boa metodologia a
pesquisa pode perder a validade. O pesquisador deve utilizar métodos ou
técnicas aceitas pela comunidade científica na busca de provar suas hipóteses.

A metodologia escolhida deve ser aquela que mais se adéqua ao seu objeto de
estudo e à abordagem aplicada. Há dois métodos principais: 1) quantitativo, que
é o uso de instrumental estatístico, de dados numéricos; e 2) qualitativo, que
se caracteriza pela qualificação dos dados coletados, durante a análise do
problema.

\section{Uma seção}
Texto.

\section{Uma outra seção}
Texto.

\chapter{Apresentação e Análise dos Resultados}

Toda pesquisa deve apresentar uma análise sobre a investigação que foi
realizada através da metodologia que foi aplicada. Nesta sessão é interessante
inserir tabelas, gráficos, imagens que mostrem os resultados, análise de dados
coletados, etc.

É interessante que nessa sessão o autor compare os seus resultados com os
resultados de outros trabalhos existentes. Essa comparação aumenta a qualidade
do trabalho e demonstra a relevância do mesmo. 

\chapter{Conclusões e Trabalhos Futuros}

A conclusão deve conter os principais aspectos e contribuições de forma a
finalizar o trabalho apresentado. Deve-se apresentar o que era esperado do
trabalho através dos objetivos inseridos inicialmente e mostrar o que foi
conseguido.

Não deve-se inserir um novo assunto na conclusão. Aqui o autor apresentará as
próprias impressões sobre o trabalho efetuado.

É importante também que sejam identificadas limitações e problemas que surgiram
durante o desenvolvimento do trabalho e quais as consequências do mesmo.

Os trabalhos futuros devem conter oportunidades de expansão do trabalho
apresentado, bem como, novos projetos que puderam ser vislumbrados a partir do
desenvolvimento do trabalho


% Referencias

\begin{references}
  \bibliography{bib/references}
\end{references}

% Apendice

\theappendix
\chapter{Mapping Study's Instruments}
\label{ap:mapping-study}

\begin{table}[!htp]
	\centering
	\caption{List of conferences on which the searches were performed.}
	\label{tbl:conferences_list}
	\rowcolors{2}{lightgray!30}{white}
	\resizebox{\columnwidth}{!}{
	\begin{tabular}{ll}
	\toprule
	\textbf{Acronym} & \textbf{Conference} \\
	\toprule
	APSEC & Asia Pacific Software Engineering Conference \\
	ASE   & IEEE/ACM International Conference on Automated Software Engineering \\
	CSMR  & European Conference on Software Maintenance and Reengineering \\
	ESEC  & European Software Engineering Conference \\
	ESEM  & International Symposium on Empirical Software Management and Measurement \\
	ICSE  & International Conference on Software Engineering \\
	ICSM  & International Conference on Software Maintenance \\
	ICST & International Conference on Software Testing \\
	InfoVis & IEEE Information Visualization Conference \\
	KDD   & ACM SIGKDD International Conference on Knowledge Discovery and Data Mining \\
	MSR   & Working Conference on Mining Software Repositories \\
	OOPSLA & Object-Oriented Programming, Systems, Languages and Applications \\
	QSIC  & International Conference On Quality Software \\
	SAC & ACM Symposium on Applied Computing \\
	SEAA & EUROMICRO Conference on Software Engineering and Advanced Applications\\
	SEDE & 19th International Conference on Software Engineering and Data Engineering \\
	SEKE  & International Conference on Software Engineering and Knowledge Engineering \\
	\bottomrule
	\end{tabular}
	}
\end{table}

\begin{table}[htp]
	\caption{List of journals in which the searches were performed.}
	\label{tbl:journals_list}
	\centering
	\rowcolors{2}{lightgray!30}{white}
	\begin{tabular}{l}
	\toprule
	\textbf{Journal title} \\
	\toprule
	ACM Transactions on Software Engineering and Methodology \\
	Automated Software Engineering \\
	Elsevier Information and Software Technology \\
	Elsevier Journal of Systems and Software \\
	Empirical Software Engineering \\
	IEEE Software \\
	IEEE Computer \\
	IEEE Transactions on Software Engineering \\
	International Journal of Software Engineering and Knowledge Engineering \\
	Journal of Software: Evolution and Process \\
	Software Quality Journal \\
	Journal of Software \\
	Software Practice and Experience Journal \\
	\bottomrule
	\end{tabular}
\end{table}

\begin{table}[h]
\centering
\footnotesize
 \rowcolors{2}{lightgray!30}{white}
\caption{Search string per Search Engine.}
\label{tbl:stringengine}
\begin{tabular}{p{.15\textwidth}p{.8\textwidth}}
\toprule
\textbf{Search Engine} & \textbf{Search String}\\
\toprule
   	 	Google Scholar &  bug report OR track OR triage ``change
   	 	request'' issue track OR request OR software OR ``modification request'' OR
   	 	``defect track'' OR ``software issue''  repositories maintenance evolution\\

   	 	ACM Portal & Abstract: "bug report" or Abstract:"change request"
   	 	or Abstract:"bug track" or Abstract:"issue track" or  Abstract:"defect
   	 	track" or Abstract:"bug triage" or Abstract: "software issue" or Abstract: "issue request"
   	 	or Abstract: "modification request") and  (Abstract:software or
   	 	Abstract:maintenance or Abstract:repositories or Abstract:repository \\

   	 	IEEExplorer (1) & (((((((((("Abstract": "bug report") OR
   	 	"Abstract":"change request") OR "Abstract":"bug track") OR "Abstract":"software issue") OR "Abstract":"issue request") OR
        "Abstract":"modification request") OR "Abstract":"issue track") OR
	    "Abstract":"defect track") OR "Abstract":"bug triage") AND
	    "Abstract":software)\\

         IEEExplorer (2) & (((((((((("Abstract": "bug report") OR
         "Abstract":"change request") OR "Abstract":"bug track") OR "Abstract":"software issue") OR
         "Abstract":"issue request") OR "Abstract":"modification request") OR
         "Abstract":"issue track") OR "Abstract":"defect track") OR
         "Abstract":"bug triage") AND "Abstract":maintenance)\\

         IEEExplorer (3) & (((((((((("Abstract": "bug report") OR
         "Abstract":"change request") OR "Abstract":"bug track") OR "Abstract":"software issue") OR
         "Abstract":"issue request") OR "Abstract":"modification request") OR
         "Abstract":"issue track") OR "Abstract":"defect track") OR
         "Abstract":"bug triage") AND "Abstract":repositories)\\

         IEEExplorer & (((((((((("Abstract": "bug report") OR
         "Abstract":"change request") OR "Abstract":"bug track") OR "Abstract":"software issue") OR
         "Abstract":"issue request") OR "Abstract":"modification request") OR
         "Abstract":"issue track") OR "Abstract":"defect track") OR
         "Abstract":"bug triage") AND "Abstract": repository)\\

         Citeseer Library & (abstract: "bug report" OR abstract:"change request" OR abstract:"bug track" OR abstract:"issue track" OR
	     abstract:"defect track" OR abstract:"bug triage" OR abstract: "software
	     issue" OR abstract: "issue request" OR abstract: "modification request")
	     AND (abstract:software OR abstract:maintenance OR abstract:repositories OR
	     abstract:repository)\\

	     Elsevier & ("bug report" OR "change
	     request" OR "bug track" OR "issue track" OR "defect track" OR "bug triage" OR "software issue" OR  "issue request" OR
	    "modification request") AND (software OR maintenance OR repositories OR
	    repository)\\

	    Scirus & ("bug report" OR "change request" OR "bug track" OR "issue track" OR  "defect track" OR "bug triage" OR
        "software issue" OR  "issue request" OR "modification request") AND
	    (software maintenance OR repositories OR repository) ANDNOT (medical OR
	    aerospace)\\

	    ScienceDirect & ("bug report" OR "change request" OR "bug track"
	     OR "issue track" OR "defect track" OR "bug triage" OR "issue request" OR
	     "modification request") AND LIMIT-TO(topics, "soft ware")\\

	     Scopus & ("bug report" OR "change request" OR "bug track" OR
	     "issue track" OR  "defect track" OR "bug triage" OR "software issue" OR
	     "issue request" OR "modification request") AND (software maintenance OR
	     repositories OR repository)\\

	     Wiley & ("bug report" OR "change request"
	     OR "bug track" OR "issue track" OR  "defect track" OR "bug triage" OR
         "software issue" OR  "issue request" OR "modification request") AND
	     (software maintenance OR repositories OR repository)\\

	     ISI Web\newline of Knowledge & ("bug report" OR "change request" OR "bug
	     track" OR "issue track" OR  "defect track" OR "bug triage" OR "software issue" OR  "issue request" OR "modification request") AND
	    (software maintenance OR repositories OR repository) ANDNOT (medical OR
	    aerospace)\\

	    SpringerLink & ("bug report" OR "change request" OR "bug track" OR "issue track" OR  "defect track" OR "bug triage" OR
        "software issue" OR  "issue request" OR "modification request") AND
	    (software maintenance OR repositories OR repository) ANDNOT (medical OR
	    aerospace)\\
	\bottomrule
\end{tabular}
\end{table}

\end{document}
